\begin{lexlean}{ResearchExecution}
\useglossary{lexlean.std.bool@1.2.0}
\useglossary{lexlean.std.nat@1.2.0}
\importmodule{LexLeanStd.Host}
\importmodule{Research}
\title{Boolean}
\begin{semanticmodule}
definition{
  name: boundary,
  parameters: [],
  result: String,
  body: "This supervisor performs one local tool job through LexLeanStd.Host. A completion is process evidence, never a Research checked result. The caller must validate the tool's report, exact contextual target, kernel replay, axiom policy and snapshot before constructing a proof answer. The driver must honor timed waits and stop ownership; a failed stop blocks further work. Before spawning, an empty zero-budget wait obtains the monotonic host clock. The second spawn result is an OS process ID, never a timestamp. Spawn latency counts against the job allowance. Cancellation drains the initial clock read before any spawn. Poll fuel also bounds a host clock that stops advancing. Tool executables resolve in the caller's pinned tool environment. No model service is a tool.",
  executable: true
}

inductive{
  name: Tool,
  type_parameters: [],
  parameters: [],
  constructors: [{name: lean, fields: []}, {name: lake, fields: []}, {name: lexlean, fields: []}, {name: checker, fields: []}, {name: git, fields: []}, {name: prismpm, fields: []}]
}

definition{
  name: program,
  parameters: [tool : Tool],
  result: String,
  body: (match tool | Tool.lean => "lean" | Tool.lake => "lake" | Tool.lexlean => "lexlean" | Tool.checker => "leanchecker" | Tool.git => "git" | Tool.prismpm => "prismpm"),
  executable: true
}

structure{
  name: Job,
  type_parameters: [],
  parameters: [],
  fields: [{name: request, type: Research::Request}, {name: tool, type: Tool}, {name: arguments, type: (List String)}, {name: directory, type: String}, {name: environment, type: (List (String * String))}, {name: output, type: String}, {name: errors, type: String}, {name: allowance, type: Nat}, {name: quantum, type: Nat}, {name: polls, type: Nat}]
}

inductive{
  name: EndKind,
  type_parameters: [],
  parameters: [],
  constructors: [{name: succeeded, fields: []}, {name: failed, fields: []}, {name: cut, fields: []}, {name: unavailable, fields: []}, {name: protocol, fields: []}, {name: cleanupFailed, fields: []}, {name: refused, fields: []}]
}

structure{
  name: Completion,
  type_parameters: [],
  parameters: [],
  fields: [{name: request, type: Research::Request}, {name: kind, type: EndKind}, {name: code, type: (Option Int32)}, {name: clock, type: Nat}, {name: output, type: String}, {name: errors, type: String}, {name: detail, type: Bytes}]
}

inductive{
  name: Phase,
  type_parameters: [],
  parameters: [],
  constructors: [{name: spawning, fields: [Nat]}, {name: cancelledSpawn, fields: [Nat]}, {name: cancelledWait, fields: [Nat, Nat]}, {name: waiting, fields: [Nat, Nat, Nat, Nat]}, {name: stopping, fields: [EndKind, Nat, Bytes]}, {name: ended, fields: [Completion]}, {name: clock, fields: []}, {name: cancelledClock, fields: []}]
}

structure{
  name: State,
  type_parameters: [],
  parameters: [],
  fields: [{name: job, type: Job}, {name: phase, type: Phase}, {name: owner, type: (Option Nat)}]
}

inductive{
  name: Emission,
  type_parameters: [],
  parameters: [],
  constructors: [{name: awaiting, fields: []}, {name: host, fields: [LexLeanStd.Host::Request]}, {name: complete, fields: [Completion]}]
}

definition{
  name: finish,
  parameters: [job : Job, kind : EndKind, code : (Option Int32), clock : Nat, detail : Bytes],
  result: (State * Emission),
  body: (let receipt : Completion := (record Completion {request := (project job request), kind := kind, code := code, clock := clock, output := (project job output), errors := (project job errors), detail := detail}); ((record State {job := job, phase := (ctor Phase.ended receipt), owner := (ctor Option.none [Nat])}), (ctor Emission.complete receipt))),
  executable: true
}

definition{
  name: validJob,
  parameters: [job : Job],
  result: Bool,
  body: ((0 <? (project job allowance)) && ((0 <? (project job quantum)) && ((0 <? (project job polls)) && ((not (prim equal (project job directory) "" : Bool)) && ((not (prim equal (project job output) "" : Bool)) && ((not (prim equal (project job errors) "" : Bool)) && (not (prim equal (project job output) (project job errors) : Bool)))))))),
  executable: true
}

definition{
  name: spawn,
  parameters: [job : Job, started : Nat],
  result: (State * Emission),
  body: ((record State {job := job, phase := (ctor Phase.spawning started), owner := (ctor Option.none [Nat])}), (ctor Emission.host (ctor LexLeanStd.Host::Request.spawn (record LexLeanStd.Host::Child {program := (program (project job tool)), arguments := (project job arguments), directory := (project job directory), environment := (project job environment), output := (project job output), errors := (project job errors)})))),
  executable: true
}

definition{
  name: begin,
  parameters: [job : Job],
  result: (State * Emission),
  body: (if (validJob job) then ((record State {job := job, phase := (ctor Phase.clock), owner := (ctor Option.none [Nat])}), (ctor Emission.host (ctor LexLeanStd.Host::Request.wait (nil Nat) 0))) else (finish job (ctor EndKind.refused) (ctor Option.none [Int32]) 0 0x696e76616c6964206a6f62)),
  axioms: [],
  executable: true,
  production: {effects: [allocation, overflow, process, recursion], targets: [rust-std]}
}

definition{
  name: waitFor,
  parameters: [job : Job, child : Nat, started : Nat, clock : Nat, polls : Nat],
  result: (State * Emission),
  body: ((record State {job := job, phase := (ctor Phase.waiting child started clock polls), owner := (ctor Option.some [Nat] child)}), (ctor Emission.host (ctor LexLeanStd.Host::Request.wait (cons child (nil Nat)) (if ((project job quantum) <? (prim subtract (project job allowance) (prim subtract clock started : Nat) : Nat)) then (project job quantum) else (prim subtract (project job allowance) (prim subtract clock started : Nat) : Nat))))),
  executable: true
}

definition{
  name: stopFor,
  parameters: [job : Job, child : Nat, clock : Nat, reason : EndKind, detail : Bytes],
  result: (State * Emission),
  body: ((record State {job := job, phase := (ctor Phase.stopping reason clock detail), owner := (ctor Option.some [Nat] child)}), (ctor Emission.host (ctor LexLeanStd.Host::Request.stop child))),
  executable: true
}

inductive{
  name: Observed,
  type_parameters: [],
  parameters: [],
  constructors: [{name: pending, fields: []}, {name: exited, fields: [(Option Int32)]}, {name: malformed, fields: []}]
}

definition{
  name: observe,
  parameters: [child : Nat, finished : (List (Nat * (Option Int32)))],
  result: Observed,
  body: (match finished
    | List.nil => (ctor Observed.pending)
    | List.cons head tail => (match tail
      | List.nil => (if ((first head) =? child) then (ctor Observed.exited (second head)) else (ctor Observed.malformed))
      | List.cons other rest => (ctor Observed.malformed))),
  executable: true
}

definition{
  name: completed,
  parameters: [job : Job, clock : Nat, code : (Option Int32)],
  result: (State * Emission),
  body: (finish job (match code | Option.none => (ctor EndKind.failed) | Option.some value => (if (prim equal value (0 : Int32) : Bool) then (ctor EndKind.succeeded) else (ctor EndKind.failed))) code clock 0x),
  executable: true
}

definition{
  name: waited,
  parameters: [job : Job, child : Nat, started : Nat, previous : Nat, polls : Nat, outcome : (Result LexLeanStd.Host::Waited LexLeanStd.Host::Failure)],
  result: (State * Emission),
  body: (match outcome
    | Result.error failure => (stopFor job child previous (ctor EndKind.unavailable) (project failure message))
    | Result.ok result => (let clock : Nat := (project result clock); (if (clock <? previous) then (stopFor job child previous (ctor EndKind.protocol) 0x636c6f636b20726567726573736564) else (match (observe child (project result finished))
      | Observed.malformed => (stopFor job child clock (ctor EndKind.protocol) 0x756e6578706563746564206368696c6420726573756c74)
      | Observed.exited code => (if ((prim subtract clock started : Nat) <? (project job allowance)) then (completed job clock code) else (finish job (ctor EndKind.cut) code clock 0x646561646c696e65))
      | Observed.pending => (if (((prim subtract clock started : Nat) <? (project job allowance)) && (1 <? polls)) then (waitFor job child started clock (prim subtract polls 1 : Nat)) else (stopFor job child clock (ctor EndKind.cut) 0x627564676574)))))),
  executable: true
}

definition{
  name: cancelledReply,
  parameters: [job : Job, child : Nat, previous : Nat, outcome : (Result LexLeanStd.Host::Waited LexLeanStd.Host::Failure)],
  result: (State * Emission),
  body: (match outcome
    | Result.error failure => (stopFor job child previous (ctor EndKind.cut) (project failure message))
    | Result.ok result => (let clock : Nat := (project result clock); (if (clock <? previous) then (stopFor job child previous (ctor EndKind.protocol) 0x636c6f636b20726567726573736564) else (match (observe child (project result finished))
      | Observed.pending => (stopFor job child clock (ctor EndKind.cut) 0x63616e63656c6c6564)
      | Observed.exited code => (finish job (ctor EndKind.cut) code clock 0x63616e63656c6c6564)
      | Observed.malformed => (stopFor job child clock (ctor EndKind.protocol) 0x756e6578706563746564206368696c6420726573756c74))))),
  executable: true
}

definition{
  name: spawned,
  parameters: [job : Job, started : Nat, outcome : (Result (Nat * Nat) LexLeanStd.Host::Failure)],
  result: (State * Emission),
  body: (match outcome
    | Result.error failure => (finish job (ctor EndKind.unavailable) (ctor Option.none [Int32]) started (project failure message))
    | Result.ok child => (waitFor job (first child) started started (project job polls))),
  executable: true
}

definition{
  name: stopped,
  parameters: [job : Job, reason : EndKind, clock : Nat, detail : Bytes, outcome : (Result Unit LexLeanStd.Host::Failure)],
  result: (State * Emission),
  body: (match outcome
    | Result.ok value => (finish job reason (ctor Option.none [Int32]) clock detail)
    | Result.error failure => (finish job (ctor EndKind.cleanupFailed) (ctor Option.none [Int32]) clock (project failure message))),
  executable: true
}

definition{
  name: asSpawned,
  parameters: [event : LexLeanStd.Host::Event],
  result: (Option (Result (Nat * Nat) LexLeanStd.Host::Failure)),
  body: (match event
    | LexLeanStd.Host::Event.spawned value => (ctor Option.some [(Result (Nat * Nat) LexLeanStd.Host::Failure)] value)
    | LexLeanStd.Host::Event.start value => (ctor Option.none [(Result (Nat * Nat) LexLeanStd.Host::Failure)])
    | LexLeanStd.Host::Event.read value => (ctor Option.none [(Result (Nat * Nat) LexLeanStd.Host::Failure)])
    | LexLeanStd.Host::Event.finished value => (ctor Option.none [(Result (Nat * Nat) LexLeanStd.Host::Failure)])
    | LexLeanStd.Host::Event.listed value => (ctor Option.none [(Result (Nat * Nat) LexLeanStd.Host::Failure)])
    | LexLeanStd.Host::Event.stat value => (ctor Option.none [(Result (Nat * Nat) LexLeanStd.Host::Failure)])
    | LexLeanStd.Host::Event.canonical value => (ctor Option.none [(Result (Nat * Nat) LexLeanStd.Host::Failure)])
    | LexLeanStd.Host::Event.waited value => (ctor Option.none [(Result (Nat * Nat) LexLeanStd.Host::Failure)])),
  executable: true
}

definition{
  name: asWaited,
  parameters: [event : LexLeanStd.Host::Event],
  result: (Option (Result LexLeanStd.Host::Waited LexLeanStd.Host::Failure)),
  body: (match event
    | LexLeanStd.Host::Event.waited value => (ctor Option.some [(Result LexLeanStd.Host::Waited LexLeanStd.Host::Failure)] value)
    | LexLeanStd.Host::Event.start value => (ctor Option.none [(Result LexLeanStd.Host::Waited LexLeanStd.Host::Failure)])
    | LexLeanStd.Host::Event.read value => (ctor Option.none [(Result LexLeanStd.Host::Waited LexLeanStd.Host::Failure)])
    | LexLeanStd.Host::Event.finished value => (ctor Option.none [(Result LexLeanStd.Host::Waited LexLeanStd.Host::Failure)])
    | LexLeanStd.Host::Event.listed value => (ctor Option.none [(Result LexLeanStd.Host::Waited LexLeanStd.Host::Failure)])
    | LexLeanStd.Host::Event.stat value => (ctor Option.none [(Result LexLeanStd.Host::Waited LexLeanStd.Host::Failure)])
    | LexLeanStd.Host::Event.canonical value => (ctor Option.none [(Result LexLeanStd.Host::Waited LexLeanStd.Host::Failure)])
    | LexLeanStd.Host::Event.spawned value => (ctor Option.none [(Result LexLeanStd.Host::Waited LexLeanStd.Host::Failure)])),
  executable: true
}

definition{
  name: asFinished,
  parameters: [event : LexLeanStd.Host::Event],
  result: (Option (Result Unit LexLeanStd.Host::Failure)),
  body: (match event
    | LexLeanStd.Host::Event.finished value => (ctor Option.some [(Result Unit LexLeanStd.Host::Failure)] value)
    | LexLeanStd.Host::Event.start value => (ctor Option.none [(Result Unit LexLeanStd.Host::Failure)])
    | LexLeanStd.Host::Event.read value => (ctor Option.none [(Result Unit LexLeanStd.Host::Failure)])
    | LexLeanStd.Host::Event.waited value => (ctor Option.none [(Result Unit LexLeanStd.Host::Failure)])
    | LexLeanStd.Host::Event.listed value => (ctor Option.none [(Result Unit LexLeanStd.Host::Failure)])
    | LexLeanStd.Host::Event.stat value => (ctor Option.none [(Result Unit LexLeanStd.Host::Failure)])
    | LexLeanStd.Host::Event.canonical value => (ctor Option.none [(Result Unit LexLeanStd.Host::Failure)])
    | LexLeanStd.Host::Event.spawned value => (ctor Option.none [(Result Unit LexLeanStd.Host::Failure)])),
  executable: true
}

definition{
  name: clocked,
  parameters: [job : Job, cancelled : Bool, event : LexLeanStd.Host::Event],
  result: (State * Emission),
  body: (match (asWaited event) | Option.none => (finish job (ctor EndKind.protocol) (ctor Option.none [Int32]) 0 0x657870656374656420636c6f636b) | Option.some outcome => (match outcome | Result.error failure => (finish job (if cancelled then (ctor EndKind.cut) else (ctor EndKind.unavailable)) (ctor Option.none [Int32]) 0 (project failure message)) | Result.ok response => (match (project response finished) | List.nil => (if cancelled then (finish job (ctor EndKind.cut) (ctor Option.none [Int32]) (project response clock) 0x63616e63656c6c6564) else (spawn job (project response clock))) | List.cons head rest => (finish job (ctor EndKind.protocol) (ctor Option.none [Int32]) (project response clock) 0x756e657870656374656420636c6f636b206368696c64)))),
  executable: true
}

definition{
  name: advance,
  parameters: [state : State, event : LexLeanStd.Host::Event],
  result: (State * Emission),
  body: (let job : Job := (project state job); (match (project state phase)
    | Phase.ended result => (state, (ctor Emission.complete result))
    | Phase.clock => (clocked job false event)
    | Phase.cancelledClock => (clocked job true event)
    | Phase.cancelledSpawn started => (match (asSpawned event)
      | Option.some outcome => (match outcome
        | Result.ok child => (stopFor job (first child) started (ctor EndKind.cut) 0x63616e63656c6c6564)
        | Result.error failure => (finish job (ctor EndKind.cut) (ctor Option.none [Int32]) 0 (project failure message)))
      | Option.none => (finish job (ctor EndKind.protocol) (ctor Option.none [Int32]) 0 0x657870656374656420737061776e))
    | Phase.cancelledWait child clock => (match (asWaited event)
      | Option.some outcome => (cancelledReply job child clock outcome)
      | Option.none => (stopFor job child clock (ctor EndKind.protocol) 0x65787065637465642077616974))
    | Phase.spawning started => (match (asSpawned event)
      | Option.some outcome => (spawned job started outcome)
      | Option.none => (finish job (ctor EndKind.protocol) (ctor Option.none [Int32]) 0 0x657870656374656420737061776e))
    | Phase.waiting child started previous polls => (match (asWaited event)
      | Option.some outcome => (waited job child started previous polls outcome)
      | Option.none => (stopFor job child previous (ctor EndKind.protocol) 0x65787065637465642077616974))
    | Phase.stopping reason clock detail => (match (asFinished event)
      | Option.some outcome => (stopped job reason clock detail outcome)
      | Option.none => (finish job (ctor EndKind.cleanupFailed) (ctor Option.none [Int32]) clock 0x65787065637465642073746f70)))),
  executable: true
}
definition{
  name: resumeAllowed,
  parameters: [completion : Completion],
  result: Bool,
  body: (match (project completion kind)
    | EndKind.succeeded => true
    | EndKind.failed => true
    | EndKind.cut => true
    | EndKind.unavailable => true
    | EndKind.protocol => false
    | EndKind.cleanupFailed => false
    | EndKind.refused => false),
  executable: true
}


definition{
  name: step,
  parameters: [state : State, event : LexLeanStd.Host::Event],
  result: (State * Emission),
  body: (let result : (State * Emission) := (advance state event); (let next : State := (first result); (let held : (Option Nat) := (match (project next owner) | Option.some child => (ctor Option.some [Nat] child) | Option.none => (project state owner)); ((record State {job := (project next job), phase := (project next phase), owner := (match (project next phase) | Phase.ended completion => (if (resumeAllowed completion) then (ctor Option.none [Nat]) else held) | Phase.clock => held | Phase.cancelledClock => held | Phase.spawning started => held | Phase.cancelledSpawn started => held | Phase.waiting child started clock polls => held | Phase.cancelledWait child clock => held | Phase.stopping kind clock detail => held)}), (second result))))),
  axioms: [],
  executable: true,
  production: {effects: [allocation, overflow, process, recursion], targets: [rust-std]}
}

definition{
  name: cancel,
  parameters: [state : State],
  result: (State * Emission),
  body: (match (project state phase)
    | Phase.clock => ((record State {job := (project state job), phase := (ctor Phase.cancelledClock), owner := (project state owner)}), (ctor Emission.awaiting))
    | Phase.cancelledClock => (state, (ctor Emission.awaiting))
    | Phase.spawning started => ((record State {job := (project state job), phase := (ctor Phase.cancelledSpawn started), owner := (project state owner)}), (ctor Emission.awaiting))
    | Phase.cancelledSpawn started => (state, (ctor Emission.awaiting))
    | Phase.waiting child started clock polls => ((record State {job := (project state job), phase := (ctor Phase.cancelledWait child clock), owner := (project state owner)}), (ctor Emission.awaiting))
    | Phase.cancelledWait child clock => (state, (ctor Emission.awaiting))
    | Phase.stopping reason clock detail => (state, (ctor Emission.awaiting))
    | Phase.ended result => (state, (ctor Emission.complete result))),
  executable: true
}


theorem{
  name: completionKeepsRequestCustody,
  parameters: [job : Job, kind : EndKind, code : (Option Int32), clock : Nat, detail : Bytes],
  statement: ((project (project (first (finish job kind code clock detail)) job) request) = (project job request)),
  proof: reflexivity{}
}

theorem{
  name: failedCleanupCannotResume,
  parameters: [request : Research::Request, code : (Option Int32), clock : Nat, output : String, errors : String, detail : Bytes],
  statement: ((resumeAllowed (record Completion {request := request, kind := (ctor EndKind.cleanupFailed), code := code, clock := clock, output := output, errors := errors, detail := detail})) = false),
  proof: reflexivity{}
}
definition{
  name: features,
  parameters: [],
  result: (List String),
  body: (cons "clock-before-spawn" (cons "process-id-is-not-time" (cons "spawn-latency-budget" (cons "bounded-timed-wait" (cons "drain-before-stop" (cons "terminal-child-ownership" (cons "cleanup-failure-retention" (nil String)))))))),
  axioms: [],
  executable: true
}

\end{semanticmodule}
\end{lexlean}
